\documentclass{article}

\usepackage{braket}
\usepackage{amsmath}
\usepackage{amssymb}

\title{Devoir 5}
\author{Thomas Lesieur et Maxime Sévigny}
\date{}
\begin{document}
\maketitle
\newpage
\section*{Question 1}
*chaque circuit pour cette question sont appelé $C_x$ où $x$ est l'ordre dans lequel ils apparaissent de droite à gauche*

\section*{1:}
Calculons le résultat du circuit $C_1$:
\begin{align*}
&t_0:&  \ket{\psi_0} &= \ket{x} \\ \\
&t_1:&  \ket{\psi_1} &= H_1\ket{\psi_0} \\
&&      \ket{\psi_1} &= H_1\ket{\psi_0} \\
&&      \ket{\psi_1} &= H_1\ket{x} \\
&&      \ket{\psi_1} &= \frac{1}{\sqrt{2}}(\ket{0} + (-1)^x\ket{1}) \\ \\
&t_2:&  \ket{\psi_2} &= X_1\ket{\psi_1} \\
&t_2:&  \ket{\psi_2} &= \frac{1}{\sqrt{2}}(X_1\ket{0} + (-1)^xX_1\ket{1}) \\
&t_2:&  \ket{\psi_2} &= \frac{1}{\sqrt{2}}(\ket{1} + (-1)^x\ket{0}) \\ \\
&t_3:&  \ket{\psi_3} &= H_1\ket{\psi_2} \\
&t_3:&  \ket{\psi_3} &= \frac{1}{\sqrt{2}}(H\ket{1} + (-1)^x\ket{0}) \\
&t_3:&  \ket{\psi_3} &= \frac{1}{\sqrt{2}}(\frac{1}{\sqrt{2}}(\ket{0} - \ket{1}) + (-1)^x\frac{1}{\sqrt{2}}(\ket{0} + \ket{1})) \\
&t_3:&  \ket{\psi_3} &= \frac{1}{2}(\ket{0} - \ket{1} + (-1)^x(\ket{0} + \ket{1})) \\
\end{align*}
\newpage
Évaluons le circuit $C_1$ pour $\ket{x} = \ket{0}$:
\begin{align*}
\ket{\psi_3} &= \frac{1}{2}(\ket{0} - \ket{1} (-1)^0(\ket{0} + \ket{1})) \\
\ket{\psi_3} &= \frac{1}{2}(\ket{0} - \ket{1} + \ket{0} + \ket{1}) \\
\ket{\psi_3} &= \frac{1}{2}(2 \cdot \ket{0}) \\
\ket{\psi_3} &= \ket{0}
\end{align*}
Évaluons le circuit $C_1$ pour $\ket{x} = \ket{0}$:
\begin{align*}
\ket{\psi_3} &= \frac{1}{2}(\ket{0} - \ket{1} (-1)^1(\ket{0} + \ket{1})) \\
\ket{\psi_3} &= \frac{1}{2}(\ket{0} - \ket{1} - \ket{0} - \ket{1}) \\
\ket{\psi_3} &= \frac{1}{2}(-2 \cdot \ket{1}) \\
\ket{\psi_3} &= -\ket{1}
\end{align*}
Le circuit donne donc le même produit que la porte $C_2$, soit la porte $Z$, c'est-à-dire : 
\[
C_1 = Z = 
    \left\{
        \begin{aligned}
            &\ket{0} \mapsto \ket{0}\\
            &\ket{1} \mapsto -\ket{1}
        \end{aligned}
    \right .
\] \\

\subsection*{2:}
Calculons le résultat du circuit $C_1$:
\begin{align*}
&t_0:&  \ket{\psi_0} &= \ket{xy} \\ \\
&t_1:&  \ket{\psi_1} &= H_1H_2\ket{\psi_0} \\
&&      \ket{\psi_1} &= H_1\ket{x}H_2\ket{y} \\
&&      \ket{\psi_1} &= \frac{1}{\sqrt{2}}(\ket{0} + (-1)^x\ket{1})\frac{1}{\sqrt{2}}(\ket{0} + (-1)^y\ket{1}) \\
&&      \ket{\psi_1} &= \frac{1}{2}(\ket{0} + (-1)^x\ket{1})(\ket{0} + (-1)^y\ket{1}) \\
&&      \ket{\psi_1} &= \frac{1}{2}(\ket{00} + (-1)^y\ket{01} + (-1)^x\ket{10} + (-1)^{x + y}\ket{11}) \\ \\
\end{align*}
\newpage
\begin{align*}
&t_2:&  \ket{\psi_2} &= CNOT_{1 \rightarrow 2}\ket{\psi_1} \\
&&      \ket{\psi_2} &= \frac{1}{2}\left(
        \begin{aligned}
            &CNOT_{1 \rightarrow 2}\ket{00} + (-1)^y      &CNOT_{1 \rightarrow 2}\ket{01} \\
     (-1)^x &CNOT_{1 \rightarrow 2}\ket{10} + (-1)^{x + y}&CNOT_{1 \rightarrow 2}\ket{11} \\
        \end{aligned}
    \right) \\
&&      \ket{\psi_2} &= \frac{1}{2}\left(
        \begin{aligned}
            &CNOT_{1 \rightarrow 2}\ket{0 \quad 0 \oplus 0} + (-1)^y      &CNOT_{1 \rightarrow 2}\ket{0 \quad 0 \oplus 1} \\
     (-1)^x &CNOT_{1 \rightarrow 2}\ket{1 \quad 1 \oplus 0} + (-1)^{x + y}&CNOT_{1 \rightarrow 2}\ket{1 \quad 1 \oplus 1} \\
        \end{aligned}
    \right) \\
&&      \ket{\psi_2} &= \frac{1}{2}(\ket{00} + (-1)^y\ket{01} + (-1)^x\ket{11} + (-1)^{x + y}\ket{10}) \\
&&      \ket{\psi_2} &= \frac{1}{2}(\ket{00} + (-1)^y\ket{01} + (-1)^{x + y}\ket{10} + (-1)^x\ket{11}) \\ \\
&t_3:&  \ket{\psi_3} &= H_1H_2\ket{\psi_2} \\
&&      \ket{\psi_3} &= \frac{1}{2}\left(
        \begin{aligned}
                 &(H\ket{0}H\ket{0}) + \\
          (-1)^y &(H\ket{0}H\ket{1}) + \\
    (-1)^{x + y} &(H\ket{1}H\ket{0}) + \\
          (-1)^x &(H\ket{1}H\ket{1}) \\
        \end{aligned}
    \right) \\
&&      \ket{\psi_3} &= \frac{1}{2}\left(
        \begin{aligned}
                 &(\frac{1}{\sqrt{2}}(\ket{0} + \ket{1})\frac{1}{\sqrt{2}}(\ket{0} + \ket{1})) + \\
          (-1)^y &(\frac{1}{\sqrt{2}}(\ket{0} + \ket{1})\frac{1}{\sqrt{2}}(\ket{0} - \ket{1})) + \\
    (-1)^{x + y} &(\frac{1}{\sqrt{2}}(\ket{0} - \ket{1})\frac{1}{\sqrt{2}}(\ket{0} + \ket{1})) + \\
          (-1)^x &(\frac{1}{\sqrt{2}}(\ket{0} - \ket{1})\frac{1}{\sqrt{2}}(\ket{0} - \ket{1})) \\
        \end{aligned}
    \right) \\
&&      \ket{\psi_3} &= \frac{1}{4}\left(
        \begin{aligned}
                 &((\ket{0} + \ket{1})(\ket{0} + \ket{1})) + \\
          (-1)^y &((\ket{0} + \ket{1})(\ket{0} - \ket{1})) + \\
    (-1)^{x + y} &((\ket{0} - \ket{1})(\ket{0} + \ket{1})) + \\
          (-1)^x &((\ket{0} - \ket{1})(\ket{0} - \ket{1})) \\
        \end{aligned}
    \right) \\
&&      \ket{\psi_3} &= \frac{1}{4}\left(
        \begin{aligned}
                 &(\ket{00} + \ket{01} + \ket{10} + \ket{11}) + \\
          (-1)^y &(\ket{00} - \ket{01} + \ket{10} - \ket{11}) + \\
    (-1)^{x + y} &(\ket{00} + \ket{01} - \ket{10} - \ket{11}) + \\
          (-1)^x &(\ket{00} - \ket{01} - \ket{10} + \ket{11}) \\
        \end{aligned}
    \right) \\
\end{align*}
\newpage
Évaluons maintenant le circuit $C_1$ pour chaque valeur de $\ket{xy}$: \\ \\
$\ket{00}:$ \\
\begin{align*}
&&      \ket{\psi_3} &= \frac{1}{4}\left(
        \begin{aligned}
                 &(\ket{00} + \ket{01} + \ket{10} + \ket{11}) + \\
          (-1)^0 &(\ket{00} - \ket{01} + \ket{10} - \ket{11}) + \\
    (-1)^{0 + 0} &(\ket{00} + \ket{01} - \ket{10} - \ket{11}) + \\
          (-1)^0 &(\ket{00} - \ket{01} - \ket{10} + \ket{11}) \\
        \end{aligned}
    \right) \\
&&      \ket{\psi_3} &= \frac{1}{4}\left(
        \begin{aligned}
                 &\ket{00} + \ket{01} + \ket{10} + \ket{11} + \\
                 &\ket{00} - \ket{01} + \ket{10} - \ket{11} + \\
                 &\ket{00} + \ket{01} - \ket{10} - \ket{11} + \\
                 &\ket{00} - \ket{01} - \ket{10} + \ket{11} \\
        \end{aligned}
    \right) \\
&&      \ket{\psi_3} &= \frac{1}{4}(4 \cdot \ket{00}) \\
&&        \ket{\psi_3} &= \ket{00}
\end{align*}
$\ket{01}:$ \\
\begin{align*}
&&      \ket{\psi_3} &= \frac{1}{4}\left(
        \begin{aligned}
                 &(\ket{00} + \ket{01} + \ket{10} + \ket{11}) + \\
          (-1)^1 &(\ket{00} - \ket{01} + \ket{10} - \ket{11}) + \\
    (-1)^{0 + 1} &(\ket{00} + \ket{01} - \ket{10} - \ket{11}) + \\
          (-1)^0 &(\ket{00} - \ket{01} - \ket{10} + \ket{11}) \\
        \end{aligned}
    \right) \\
&&      \ket{\psi_3} &= \frac{1}{4}\left(
        \begin{aligned}
                 &\ket{00} + \ket{01} + \ket{10} + \ket{11} + \\
                -&\ket{00} + \ket{01} - \ket{10} + \ket{11} + \\
                -&\ket{00} - \ket{01} + \ket{10} + \ket{11} + \\
                 &\ket{00} - \ket{01} - \ket{10} + \ket{11} \\
        \end{aligned}
    \right) \\
&&      \ket{\psi_3} &= \frac{1}{4}(4 \cdot \ket{11}) \\
&&      \ket{\psi_3} &= \ket{11} \\
\end{align*}
\newpage
$\ket{10}:$ \\
\begin{align*}
&&      \ket{\psi_3} &= \frac{1}{4}\left(
        \begin{aligned}
                 &(\ket{00} + \ket{01} + \ket{10} + \ket{11}) + \\
          (-1)^0 &(\ket{00} - \ket{01} + \ket{10} - \ket{11}) + \\
    (-1)^{1 + 0} &(\ket{00} + \ket{01} - \ket{10} - \ket{11}) + \\
          (-1)^1 &(\ket{00} - \ket{01} - \ket{10} + \ket{11}) \\
        \end{aligned}
    \right) \\
&&      \ket{\psi_3} &= \frac{1}{4}\left(
        \begin{aligned}
                 &\ket{00} + \ket{01} + \ket{10} + \ket{11} + \\
                 &\ket{00} - \ket{01} + \ket{10} - \ket{11} + \\
                -&\ket{00} - \ket{01} + \ket{10} + \ket{11} + \\
                -&\ket{00} + \ket{01} + \ket{10} - \ket{11} \\
        \end{aligned}
    \right) \\
&&      \ket{\psi_3} &= \frac{1}{4}(4 \cdot \ket{10}) \\
&&        \ket{\psi_3} &= \ket{10}
\end{align*}
$\ket{11}:$ \\
\begin{align*}
&&      \ket{\psi_3} &= \frac{1}{4}\left(
        \begin{aligned}
                 &(\ket{00} + \ket{01} + \ket{10} + \ket{11}) + \\
          (-1)^1 &(\ket{00} - \ket{01} + \ket{10} - \ket{11}) + \\
    (-1)^{1 + 1} &(\ket{00} + \ket{01} - \ket{10} - \ket{11}) + \\
          (-1)^1 &(\ket{00} - \ket{01} - \ket{10} + \ket{11}) \\
        \end{aligned}
    \right) \\
&&      \ket{\psi_3} &= \frac{1}{4}\left(
        \begin{aligned}
                 &\ket{00} + \ket{01} + \ket{10} + \ket{11} + \\
                -&\ket{00} + \ket{01} - \ket{10} + \ket{11} + \\
                 &\ket{00} + \ket{01} - \ket{10} - \ket{11} + \\
                -&\ket{00} + \ket{01} + \ket{10} - \ket{11} \\
        \end{aligned}
    \right) \\
&&      \ket{\psi_3} &= \frac{1}{4}(4 \cdot \ket{01}) \\
&&      \ket{\psi_3} &= \ket{01} \\
\end{align*}
Le circuit $C_2$ correspond à l'équation suivante : \\
\begin{align*} CNOT_{2 \rightarrow 1}\ket{xy} = \ket{x \oplus y \quad y} \end{align*}
\newpage
Évaluons $C_2$ pour chaque $\ket{xy}$ possible:
\begin{align*}
CNOT_{2 \rightarrow 1}\ket{00} &= \ket{0 \oplus 0 \quad 0} = \ket{00} \\
CNOT_{2 \rightarrow 1}\ket{01} &= \ket{0 \oplus 1 \quad 1} = \ket{11} \\ 
CNOT_{2 \rightarrow 1}\ket{10} &= \ket{1 \oplus 0 \quad 0} = \ket{10} \\ 
CNOT_{2 \rightarrow 1}\ket{11} &= \ket{1 \oplus 1 \quad 1} = \ket{01} \\
\end{align*}
Les valeurs obtenus pour les circuits $C_1$ et $C_2$ sont donc identiques pour tout $x \in \{0,1\}, y \in \{0, 1\}$: \\
\begin{align*}
C_1\ket{00} &= C_2 \ket{00} &= \ket{00} \\
C_1\ket{01} &= C_2 \ket{01} &= \ket{11} \\
C_1\ket{10} &= C_2 \ket{10} &= \ket{10} \\ 
C_1\ket{11} &= C_2 \ket{11} &= \ket{01} \\
\end{align*}

\subsection*{3:}
Calculons le résultat du circuit $C_1$:
\begin{align*}
&t_0:&  \ket{\psi_0} &= \ket{xy} \\ \\
&t_1:&  \ket{\psi_1} &= H_2\ket{\psi_0} \\
&&      \ket{\psi_1} &= \ket{x}H\ket{y} \\
&&      \ket{\psi_1} &= \ket{x}\frac{1}{\sqrt{2}}(\ket{0} + (-1)^y \ket{1}) \\
&&      \ket{\psi_1} &= \frac{1}{\sqrt{2}}(\ket{x0} + (-1)^y \ket{x1}) \\ \\
&t_2:&  \ket{\psi_2} &= CNOT_{1 \rightarrow 2}\ket{\psi_1} \\
&&      \ket{\psi_2} &= \frac{1}{\sqrt{2}}(CNOT_{1 \rightarrow 2}\ket{x0} + (-1)^yCNOT_{1 \rightarrow 2}\ket{x1}) \\
&&      \ket{\psi_2} &= \frac{1}{\sqrt{2}}(\ket{x \quad x \oplus 0} + (-1)^y\ket{x \quad x \oplus 1}) \\
&&      \ket{\psi_2} &= \frac{1}{\sqrt{2}}(\ket{xx} + (-1)^y\ket{x \overline{x}}) \\ \\
\end{align*}
\begin{align*}
&t_3:&  \ket{\psi_3} &= H_2\ket{\psi_2} \\
&&      \ket{\psi_3} &= \frac{1}{\sqrt{2}}(\ket{x}H\ket{x} + (-1)^y \ket{x}H\ket{\overline{x}}) \\
&&      \ket{\psi_3} &= \frac{1}{\sqrt{2}}\left(
        \begin{aligned}
                &\ket{x}\frac{1}{\sqrt{2}}(\ket{0} + (-1)^x \ket{1}) + \\
        (-1)^y  &\ket{x}\frac{1}{\sqrt{2}}(\ket{0} + (-1)^{\overline{x}} \ket{1}) \\
        \end{aligned}
    \right) \\
&&      \ket{\psi_3} &= \frac{1}{2}\left(
        \begin{aligned}
                &\ket{x}(\ket{0} + (-1)^x \ket{1}) + \\
        (-1)^y  &\ket{x}(\ket{0} + (-1)^{\overline{x}} \ket{1}) \\
        \end{aligned}
    \right) \\
&&      \ket{\psi_3} &= \frac{1}{2}(\ket{x0} + (-1)^x\ket{x1} + (-1)^y\ket{x0} + (-1)^y(-1)^{\overline{x}} \ket{x1}) \\
\end{align*}
Calculons le résultat du circuit $C_2$:
\begin{align*}
&t_0:&  \ket{\psi_0} &= \ket{xy} \\ \\
&t_1:&  \ket{\psi_1} &= H_2\ket{\psi_0} \\
&&      \ket{\psi_1} &= H\ket{x}\ket{y} \\
&&      \ket{\psi_1} &= \frac{1}{\sqrt{2}}(\ket{0} + (-1)^x\ket{1})\ket{y} \\
&&      \ket{\psi_1} &= \frac{1}{\sqrt{2}}(\ket{0y} + (-1)^x \ket{1y}) \\ \\
&t_2:&  \ket{\psi_2} &= CNOT_{2 \rightarrow 1}\ket{\psi_1} \\
&&      \ket{\psi_2} &= \frac{1}{\sqrt{2}}(CNOT_{2 \rightarrow 1}\ket{0y} + (-1)^xCNOT_{2 \rightarrow 1}\ket{1y}) \\
&&      \ket{\psi_2} &= \frac{1}{\sqrt{2}}(\ket{0 \oplus y \quad y} + (-1)^x\ket{1 \oplus y \quad y}) \\
&&      \ket{\psi_2} &= \frac{1}{\sqrt{2}}(\ket{yy} + (-1)^x\ket{\overline{y}y}) \\ \\
\end{align*}
\begin{align*}
&t_3:&  \ket{\psi_3} &= H_1\ket{\psi_2} \\
&&      \ket{\psi_3} &= \frac{1}{\sqrt{2}}(H\ket{y}\ket{y} + (-1)^xH\ket{\overline{y}}\ket{y}) \\
&&      \ket{\psi_3} &= \frac{1}{\sqrt{2}}\left(
        \begin{aligned}
                &\frac{1}{\sqrt{2}}(\ket{0} + (-1)^y\ket{1})\ket{y} + \\
        (-1)^x  &\frac{1}{\sqrt{2}}(\ket{0} + (-1)^{\overline{y}}\ket{1})\ket{y} \\
        \end{aligned}
    \right) \\
&&      \ket{\psi_3} &= \frac{1}{2}\left(
        \begin{aligned}
                &(\ket{0} + (-1)^y\ket{1})\ket{y} + \\
        (-1)^x  &(\ket{0} + (-1)^{\overline{y}}\ket{1})\ket{y} \\
        \end{aligned}
    \right) \\
&&      \ket{\psi_3} &= \frac{1}{2}(\ket{0y} + (-1)^y\ket{1y} + (-1)^x\ket{0y} + (-1)^x(-1)^{\overline{y}}\ket{1y}) \\
\end{align*}
L'évaluation des circuits $C_3$ et $C_4$ est la suivante:
\[
X_{1 \rightarrow 2} = X_{2 \rightarrow 1} = 
    \left\{
        \begin{aligned}
            &\ket{00} \mapsto&  \ket{00} \\
            &\ket{01} \mapsto&  \ket{01} \\
            &\ket{10} \mapsto&  \ket{10} \\
            &\ket{11} \mapsto& -\ket{11} \\
        \end{aligned}
    \right .
\] \\
Évaluons maintenant le circuit $C_1$ pour chaque valeur de $\ket{xy}$: \\ \\
$\ket{00}:$ \\
\begin{align*}
&&      \ket{\psi_3} &= \frac{1}{2}(\ket{00} + (-1)^0\ket{01} + (-1)^0\ket{00} + (-1)^0(-1)^1 \ket{01}) \\
&&      \ket{\psi_3} &= \frac{1}{2}(\ket{00} + \ket{01} + \ket{00} - \ket{01}) \\
&&      \ket{\psi_3} &= \frac{1}{2}(2 \cdot \ket{00}) \\
&&      \ket{\psi_3} &= \ket{00} \\
\end{align*}
$\ket{01}:$ \\
\begin{align*}
&&      \ket{\psi_3} &= \frac{1}{2}(\ket{00} + (-1)^0\ket{01} + (-1)^1\ket{00} + (-1)^1(-1)^1 \ket{01}) \\
&&      \ket{\psi_3} &= \frac{1}{2}(\ket{00} + \ket{01} - \ket{00} + \ket{01}) \\
&&      \ket{\psi_3} &= \frac{1}{2}(2 \cdot \ket{01}) \\
&&      \ket{\psi_3} &= \ket{01} \\
\end{align*}
$\ket{10}:$ \\
\begin{align*}
&&      \ket{\psi_3} &= \frac{1}{2}(\ket{10} + (-1)^1\ket{11} + (-1)^0\ket{10} + (-1)^0(-1)^0 \ket{11}) \\
&&      \ket{\psi_3} &= \frac{1}{2}(\ket{10} - \ket{11} + \ket{10} + \ket{11}) \\
&&      \ket{\psi_3} &= \frac{1}{2}(2 \cdot \ket{10}) \\
&&      \ket{\psi_3} &= \ket{10} \\
\end{align*}
$\ket{11}:$ \\
\begin{align*}
&&      \ket{\psi_3} &= \frac{1}{2}(\ket{10} + (-1)^1\ket{11} + (-1)^1\ket{10} + (-1)^1(-1)^{0} \ket{11}) \\
&&      \ket{\psi_3} &= \frac{1}{2}(\ket{10} - \ket{11} - \ket{00} - \ket{11}) \\
&&      \ket{\psi_3} &= \frac{1}{2}(-2 \cdot \ket{11}) \\
&&      \ket{\psi_3} &= -\ket{11} \\
\end{align*}
Évaluons ensuite le circuit $C_2$ pour chaque valeur de $\ket{xy}$: \\ \\
$\ket{00}:$ \\
\begin{align*}
&&      \ket{\psi_3} &= \frac{1}{2}(\ket{00} + (-1)^0\ket{10} + (-1)^0\ket{00} + (-1)^0(-1)^1\ket{10}) \\
&&      \ket{\psi_3} &= \frac{1}{2}(\ket{00} + \ket{10} + \ket{00} - \ket{10}) \\
&&      \ket{\psi_3} &= \frac{1}{2}(2 \cdot \ket{00}) \\
&&      \ket{\psi_3} &= \ket{00} \\
\end{align*}
$\ket{01}:$ \\
\begin{align*}
&&      \ket{\psi_3} &= \frac{1}{2}(\ket{01} + (-1)^1\ket{11} + (-1)^0\ket{01} + (-1)^0(-1)^0\ket{11}) \\
&&      \ket{\psi_3} &= \frac{1}{2}(\ket{01} - \ket{11} + \ket{01} + \ket{11}) \\
&&      \ket{\psi_3} &= \frac{1}{2}(2 \cdot \ket{01}) \\
&&      \ket{\psi_3} &= \ket{01} \\
\end{align*}
$\ket{10}:$ \\
\begin{align*}
&&      \ket{\psi_3} &= \frac{1}{2}(\ket{00} + (-1)^0\ket{10} + (-1)^1\ket{00} + (-1)^1(-1)^1\ket{10}) \\
&&      \ket{\psi_3} &= \frac{1}{2}(\ket{10} - \ket{11} + \ket{10} + \ket{11}) \\
&&      \ket{\psi_3} &= \frac{1}{2}(2 \cdot \ket{10}) \\
&&      \ket{\psi_3} &= \ket{10} \\
\end{align*}
$\ket{11}:$ \\
\begin{align*}
&&      \ket{\psi_3} &= \frac{1}{2}(\ket{01} + (-1)^1\ket{11} + (-1)^1\ket{01} + (-1)^1(-1)^0\ket{11}) \\
&&      \ket{\psi_3} &= \frac{1}{2}(\ket{01} - \ket{11} - \ket{01} - \ket{11}) \\
&&      \ket{\psi_3} &= \frac{1}{2}(-2 \cdot \ket{11}) \\
&&      \ket{\psi_3} &= -\ket{11} \\
\end{align*}
On Remarque donc que les circuits donnent les mêmes résultat pour tous les $\ket{xy}$ possibles:
\[
C_1\ket{xy} = C_2\ket{xy} = C_3\ket{xy} = C_4\ket{xy} =
    \left\{
        \begin{aligned}
           -&\ket{xy} si\; x = 1, y = 1 \\
            &\ket{xy} sinon \\
        \end{aligned}
    \right .
\] \\

\newpage
\subsection*{4:}
Calculons le résultat du circuit $C_1$:
\begin{align*}
&t_0:&  \ket{\psi_0} &= \ket{zx} \\ \\
&t_1:&  \ket{\psi_1} &= H_1\ket{\psi_0} \\
&&      \ket{\psi_1} &=  H\ket{z}\ket{x} \\
&&      \ket{\psi_1} &= \frac{1}{\sqrt{2}}(\ket{0} + (-1)^z\ket{1})\ket{x} \\
&&      \ket{\psi_1} &= \frac{1}{\sqrt{2}}(\ket{0x} + (-1)^z\ket{1x}) \\ \\
&t_2:&  \ket{\psi_2} &= CNOT_{1 \rightarrow 2}\ket{\psi_1} \\
&&      \ket{\psi_2} &= \frac{1}{\sqrt{2}}(CNOT_{1 \rightarrow 2}\ket{0x} + (-1)^zCNOT_{1 \rightarrow 2}\ket{1x}) \\
&&      \ket{\psi_2} &= \frac{1}{\sqrt{2}}(\ket{0 \quad 0 \oplus x} + (-1)^z\ket{1 \quad 1 \oplus x}) \\
&&      \ket{\psi_2} &= \frac{1}{\sqrt{2}}(\ket{0x} + (-1)^z\ket{1\overline{x}}) \\
\end{align*}
Évaluons maintenant le circuit $C_1$ pour chaque valeur de $\ket{zx}$: \\ \\
$\ket{00}:$ \\
\begin{align*}
&&      \ket{\psi_2} &= \frac{1}{\sqrt{2}}(\ket{00} + (-1)^0\ket{11}) \\
&&      \ket{\psi_2} &= \frac{1}{\sqrt{2}}(\ket{00} + \ket{11})\\
\end{align*}
$\ket{01}:$ \\
\begin{align*}
&&      \ket{\psi_2} &= \frac{1}{\sqrt{2}}(\ket{01} + (-1)^0\ket{10}) \\
&&      \ket{\psi_2} &= \frac{1}{\sqrt{2}}(\ket{01} + \ket{10})\\
\end{align*}
$\ket{10}:$ \\
\begin{align*}
&&      \ket{\psi_2} &= \frac{1}{\sqrt{2}}(\ket{00} + (-1)^1\ket{11}) \\
&&      \ket{\psi_2} &= \frac{1}{\sqrt{2}}(\ket{00} - \ket{11})\\
\end{align*}
$\ket{11}:$ \\
\begin{align*}
&&      \ket{\psi_2} &= \frac{1}{\sqrt{2}}(\ket{01} + (-1)^0\ket{10}) \\
&&      \ket{\psi_2} &= \frac{1}{\sqrt{2}}(\ket{01} - \ket{10})\\
\end{align*}
Calculons le résultat du circuit $C_2$:
\begin{align*}
&t_0:&  \ket{\psi_0} &= \ket{00} \\ \\
&t_1:&  \ket{\psi_1} &= H_1\ket{\psi_0} \\
&&      \ket{\psi_1} &= H\ket{0}\ket{0} \\
&&      \ket{\psi_1} &= \frac{1}{\sqrt{2}}(\ket{0} + \ket{1})\ket{0} \\
&&      \ket{\psi_1} &= \frac{1}{\sqrt{2}}(\ket{00} + \ket{10}) \\ \\
&t_2:&  \ket{\psi_2} &= CNOT_{1 \rightarrow 2}\ket{\psi_1} \\
&&      \ket{\psi_2} &= \frac{1}{\sqrt{2}}(CNOT_{1 \rightarrow 2}\ket{00} + CNOT_{1 \rightarrow 2}\ket{10}) \\
&&      \ket{\psi_2} &= \frac{1}{\sqrt{2}}(\ket{0 \quad 0 \oplus 0} + \ket{1 \quad 1 \oplus 0}) \\
&&      \ket{\psi_2} &= \frac{1}{\sqrt{2}}(\ket{00} + \ket{11}) \\ \\
\end{align*}
\begin{align*}
&t_3:&  \ket{\psi_3} &= X^x\ket{\psi_2} \\
&&      \ket{\psi_3} &= \frac{1}{\sqrt{2}}(X^x\ket{0}\ket{0} + X^x\ket{1}\ket{1}) \\
&&      \ket{\psi_3} &= \frac{1}{\sqrt{2}}(\ket{x \oplus 0 \quad 0} + \ket{x \oplus 1 \quad 1}) \\
&&      \ket{\psi_3} &= \frac{1}{\sqrt{2}}(\ket{x0} + \ket{\overline{x}1}) \\ \\
&t_4:&  \ket{\psi_4} &=  Z^z\ket{\psi_3} \\
&&      \ket{\psi_4} &= \frac{1}{\sqrt{2}}(Z^z\ket{x}\ket{0} + Z^z\ket{\overline{x}1}) \\
&&      \ket{\psi_4} &= \frac{1}{\sqrt{2}}((-1)^{z \wedge x}\ket{x}\ket{0} + (-1)^{z \wedge \overline{x}}\ket{\overline{x}}\ket{1}) \\
&&      \ket{\psi_4} &= \frac{1}{\sqrt{2}}((-1)^{z \wedge x}\ket{x0} + (-1)^{z \wedge \overline{x}}\ket{\overline{x}1}) \\
\end{align*}
Évaluons maintenant le circuit $C_2$ pour chaque valeur de $\ket{zx}$: \\ \\
$\ket{00}:$ \\
\begin{align*}
&&      \ket{\psi_2} &= \frac{1}{\sqrt{2}}((-1)^{0 \wedge 0}\ket{00} + (-1)^{0 \wedge 1}\ket{11}) \\
&&      \ket{\psi_2} &= \frac{1}{\sqrt{2}}(\ket{00} + \ket{11})\\
\end{align*}
$\ket{01}:$ \\
\begin{align*}
&&      \ket{\psi_2} &= \frac{1}{\sqrt{2}}((-1)^{0 \wedge 1}\ket{10} + (-1)^{0 \wedge 0}\ket{01}) \\
&&      \ket{\psi_2} &= \frac{1}{\sqrt{2}}(\ket{01} + \ket{10})\\
\end{align*}
$\ket{10}:$ \\
\begin{align*}
&&      \ket{\psi_2} &= \frac{1}{\sqrt{2}}((-1)^{1 \wedge 0}\ket{00} + (-1)^{1 \wedge 1}\ket{11}) \\
&&      \ket{\psi_2} &= \frac{1}{\sqrt{2}}(\ket{00} - \ket{11})\\
\end{align*}
$\ket{11}:$ \\
\begin{align*}
&&      \ket{\psi_2} &= \frac{1}{\sqrt{2}}((-1)^{1 \wedge 1}\ket{10} + (-1)^{1 \wedge 0}\ket{01}) \\
&&      \ket{\psi_2} &= \frac{1}{\sqrt{2}}(\ket{01} - \ket{10})\\
\end{align*}
Il est donc possible de déterminer que $C_1$ et $C_2$ sont des circuits identiques qui génèrent: \\
\[
C_1\ket{zx} = C_2\ket{zx} =
    \left\{
        \begin{aligned}
           \ket{00} \mapsto&\quad \frac{1}{\sqrt{2}}(\ket{00} + \ket{11}) \\
           \ket{01} \mapsto&\quad \frac{1}{\sqrt{2}}(\ket{01} + \ket{10}) \\
           \ket{10} \mapsto&\quad \frac{1}{\sqrt{2}}(\ket{00} - \ket{11}) \\
           \ket{11} \mapsto&\quad \frac{1}{\sqrt{2}}(\ket{01} - \ket{10}) \\
        \end{aligned}
    \right .
\] \\

\subsection*{5:}
Calculons le résultat du circuit $C_1$:
\begin{align*}
&t_0:&  \ket{\psi_0} &= \ket{xy} \\ \\
&t_1:&  \ket{\psi_1} &= H_1\ket{\psi_0} \\
&&      \ket{\psi_1} &= H\ket{x}\ket{y} \\
&&      \ket{\psi_1} &= \frac{1}{\sqrt{2}}(\ket{0} + (-1)^x\ket{1})\ket{y} \\
&&      \ket{\psi_1} &= \frac{1}{\sqrt{2}}(\ket{0y} + (-1)^x\ket{1y}) \\ \\
&t_2:&  \ket{\psi_2} &= CNOT_{1 \rightarrow 2}\ket{\psi_1} \\
&&      \ket{\psi_2} &= \frac{1}{\sqrt{2}}(CNOT_{1 \rightarrow 2}\ket{0y} + (-1)^xCNOT_{1 \rightarrow 2}\ket{1y}) \\
&&      \ket{\psi_2} &= \frac{1}{\sqrt{2}}(\ket{0 \quad 0 \oplus y} + (-1)^x\ket{1 \quad 1 \oplus y}) \\
&&      \ket{\psi_2} &= \frac{1}{\sqrt{2}}(\ket{0y} + (-1)^x\ket{1\overline{y}}) \\
\end{align*}
\begin{align*}
&t_3:&  \ket{\psi_3} &= CNOT_{1 \rightarrow 2}\ket{\psi_2} \\
&&      \ket{\psi_3} &= \frac{1}{\sqrt{2}}(CNOT_{1 \rightarrow 2}\ket{0y} + (-1)^xCNOT_{1 \rightarrow 2}\ket{1\overline{y}}) \\
&&      \ket{\psi_3} &= \frac{1}{\sqrt{2}}(\ket{0 \quad 0 \oplus y} + (-1)^x\ket{1 \quad 1 \oplus \overline{y}}) \\
&&      \ket{\psi_3} &= \frac{1}{\sqrt{2}}(\ket{0y} + (-1)^x\ket{1y}) \\ \\
&t_4:&  \ket{\psi_4} &= H_1\ket{\psi_3} \\
&&      \ket{\psi_4} &= \frac{1}{\sqrt{2}}(H\ket{0}\ket{y} + (-1)^xH\ket{1}\ket{y}) \\
&&      \ket{\psi_4} &= \frac{1}{\sqrt{2}}(\frac{1}{\sqrt{2}}(\ket{0} + \ket{1})\ket{y} + (-1)^x\frac{1}{\sqrt{2}}(\ket{0} - \ket{1})\ket{y}) \\
&&      \ket{\psi_4} &= \frac{1}{2}((\ket{0} + \ket{1})\ket{y} + (-1)^x(\ket{0} - \ket{1})\ket{y}) \\
&&      \ket{\psi_4} &= \frac{1}{2}(\ket{0y} + \ket{1y} + (-1)^x(\ket{0y} - \ket{1y})) \\
\end{align*}
Évaluons maintenant le circuit $C_1$ pour chaque valeur de $\ket{xy}$: \\ \\
$\ket{00}:$ \\
\begin{align*}
&&      \ket{\psi_4} &= \frac{1}{2}(\ket{00} + \ket{10} + \ket{00} - \ket{10}) \\
&&      \ket{\psi_4} &= \frac{1}{2}(2 \cdot \ket{00}) \\
&&      \ket{\psi_4} &= \ket{00} \\
\end{align*}
$\ket{01}:$ \\
\begin{align*}
&&      \ket{\psi_4} &= \frac{1}{2}(\ket{01} + \ket{11} + \ket{01} - \ket{11}) \\
&&      \ket{\psi_4} &= \frac{1}{2}(2 \cdot \ket{01}) \\
&&      \ket{\psi_4} &= \ket{01} \\
\end{align*}
$\ket{10}:$ \\
\begin{align*}
&&      \ket{\psi_4} &= \frac{1}{2}(\ket{00} + \ket{10} - \ket{00} + \ket{10}) \\
&&      \ket{\psi_4} &= \frac{1}{2}(2 \cdot \ket{10}) \\
&&      \ket{\psi_4} &= \ket{10} \\
\end{align*}
$\ket{11}:$ \\
\begin{align*}
&&      \ket{\psi_4} &= \frac{1}{2}(\ket{01} + \ket{11} - \ket{01} + \ket{11}) \\
&&      \ket{\psi_4} &= \frac{1}{2}(2 \cdot \ket{11}) \\
&&      \ket{\psi_4} &= \ket{11} \\
\end{align*}
Il est donc possible de déterminer que $C_1$ et $C_2$ sont des circuits identiques qui génèrent: \\
\[
C_1\ket{xy} = C_2\ket{xy} = \ket{xy} \mapsto \ket{xy}
\]

\section*{Question 2}
\subsection*{(a)}

Voici le développement de la fonction $f$ sur $\hat{x} = 0^2$
\begin{align*}
&t_0:& \ket{\psi_0} &= \ket{00}\ket{-} \\ \\
&t_1:& \ket{\psi_1} &= H_1H_2\ket{\psi_0} \\
&&     \ket{\psi_1} &= H_1\ket{0}H_2\ket{0}\ket{-} \\
&&     \ket{\psi_1} &= \frac{1}{\sqrt{2}}(\ket{0} + \ket{1})\frac{1}{\sqrt{2}}(\ket{0} + \ket{1})\ket{-} \\
&&     \ket{\psi_1} &= \frac{1}{2\sqrt{2}}(\ket{0} + \ket{1})(\ket{0} + \ket{1})(\ket{0} - \ket{1}) \\
&&     \ket{\psi_1} &=\frac{1}{2\sqrt{2}}(\ket{00} + \ket{01} + \ket{10} + \ket{11})(\ket{0} - \ket{1}) \\
&&     \ket{\psi_1}&= \frac{1}{2\sqrt{2}}(\ket{000} - \ket{001} + \ket{010} - \ket{011} + \ket{100} - \ket{101} + \ket{110} - \ket{111}) \\ \\
&t_2:& \ket{\psi_2} &= f_{12 \rightarrow 3}\ket{\psi_1} \\
&&     \ket{\psi_2} &= \frac{1}{2\sqrt{2}}\left(
        \begin{aligned}
            &f_{12 \rightarrow 3}\ket{000} + f_{12 \rightarrow 3}\ket{001} + \\
            &f_{12 \rightarrow 3}\ket{010} + f_{12 \rightarrow 3}\ket{011} + \\
            &f_{12 \rightarrow 3}\ket{100} + f_{12 \rightarrow 3}\ket{101} + \\
            &f_{12 \rightarrow 3}\ket{110} + f_{12 \rightarrow 3}\ket{111} \\
        \end{aligned}
    \right) \\
&&     \ket{\psi_2} &= \frac{1}{2\sqrt{2}}\left(
        \begin{aligned}
            &\ket{00}\ket{f(00) \oplus 0} - \ket{00}\ket{f(00) \oplus 1} + \\
            &\ket{01}\ket{f(01) \oplus 0} - \ket{01}\ket{f(01) \oplus 1} + \\
            &\ket{10}\ket{f(10) \oplus 0} - \ket{10}\ket{f(10) \oplus 1} + \\
            &\ket{11}\ket{f(11) \oplus 0} - \ket{11}\ket{f(11) \oplus 1}   \\
        \end{aligned}
    \right) \\
&&     \ket{\psi_2} &= \frac{1}{2\sqrt{2}}\left(
        \begin{aligned}
            &\ket{00}\ket{1 \oplus 0} - \ket{00}\ket{1 \oplus 1} + \\
            &\ket{01}\ket{0 \oplus 0} - \ket{01}\ket{0 \oplus 1} + \\
            &\ket{10}\ket{0 \oplus 0} - \ket{10}\ket{0 \oplus 1} + \\
            &\ket{11}\ket{0 \oplus 0} - \ket{11}\ket{0 \oplus 1}   \\
        \end{aligned}
    \right) \\
\end{align*}

\begin{align*}
&&     \ket{\psi_2} &= \frac{1}{2\sqrt{2}}\left(
        \begin{aligned}
            &\ket{001} - \ket{000} + \\
            &\ket{010} - \ket{011} + \\
            &\ket{100} - \ket{101} + \\
            &\ket{110} - \ket{111}   \\
        \end{aligned}
    \right) \\
&&     \ket{\psi_2} &= \frac{1}{2\sqrt{2}}(\ket{001} - \ket{000} + \ket{010} - \ket{011} + \ket{100} - \ket{101} + \ket{110} - \ket{111}) \\ \\
&t_3:& \ket{\psi_3} &= H_1H_2\ket{\psi_2} \\
&&     \ket{\psi_3} &= \frac{1}{2\sqrt{2}}\left(
        \begin{aligned}
            &H\ket{0}H\ket{0}\ket{1} - H\ket{0}H\ket{0}\ket{0} + \\
            &H\ket{0}H\ket{1}\ket{0} - H\ket{0}H\ket{1}\ket{1} + \\
            &H\ket{1}H\ket{0}\ket{0} - H\ket{1}H\ket{0}\ket{1} + \\
            &H\ket{1}H\ket{1}\ket{0} - H\ket{1}H\ket{1}\ket{1}   \\
        \end{aligned}
    \right) \\
&&  \ket{\psi_3} &= \frac{1}{4\sqrt{2}}\left(
    \begin{aligned}
        &(\ket{0} + \ket{1})(\ket{0} + \ket{1})\ket{1} - (\ket{0} + \ket{1})(\ket{0} + \ket{1})\ket{0} + \\
        &(\ket{0} + \ket{1})(\ket{0} - \ket{1})\ket{0} - (\ket{0} + \ket{1})(\ket{0} - \ket{1})\ket{1} + \\
        &(\ket{0} - \ket{1})(\ket{0} + \ket{1})\ket{0} - (\ket{0} - \ket{1})(\ket{0} + \ket{1})\ket{1} + \\
        &(\ket{0} - \ket{1})(\ket{0} - \ket{1})\ket{0} - (\ket{0} - \ket{1})(\ket{0} - \ket{1})\ket{1}   \\
    \end{aligned}
\right) \\
&&  \ket{\psi_3} &= \frac{1}{4\sqrt{2}}\left(
    \begin{aligned}
        &(\ket{0} + \ket{1})(\ket{0} + \ket{1})(\ket{1} - \ket{0}) + \\
        &(\ket{0} + \ket{1})(\ket{0} - \ket{1})(\ket{0} - \ket{1}) + \\
        &(\ket{0} - \ket{1})(\ket{0} + \ket{1})(\ket{0} - \ket{1}) + \\
        &(\ket{0} - \ket{1})(\ket{0} - \ket{1})(\ket{0} - \ket{1})   \\
    \end{aligned}
\right) \\
&&     \ket{\psi_3} &= \frac{1}{4\sqrt{2}}\left(
    \begin{aligned}
        &(\ket{00} + \ket{01} + \ket{10} + \ket{11})(\ket{1} - \ket{0}) + \\
        &(\ket{00} - \ket{01} + \ket{10} - \ket{11})(\ket{0} - \ket{1}) + \\
        &(\ket{00} + \ket{01} - \ket{10} - \ket{11})(\ket{0} - \ket{1}) + \\
        &(\ket{00} - \ket{01} - \ket{10} + \ket{11})(\ket{0} - \ket{1})   \\
    \end{aligned}
\right) \\
&&     \ket{\psi_3} &= \frac{1}{4\sqrt{2}}\left(
    \begin{aligned}
       -&\ket{000} + \ket{001} - \ket{010} + \ket{011} - \ket{100} + \ket{101} - \ket{110} + \ket{111} + \\
        &\ket{000} - \ket{001} - \ket{010} + \ket{011} + \ket{100} - \ket{101} - \ket{110} + \ket{111} + \\
        &\ket{000} - \ket{001} + \ket{010} - \ket{011} - \ket{100} + \ket{101} - \ket{110} + \ket{111} + \\
        &\ket{000} - \ket{001} - \ket{010} + \ket{011} - \ket{100} + \ket{101} + \ket{110} - \ket{111}   \\
    \end{aligned}
\right) \\
&&     \ket{\psi_3} &= \frac{1}{4\sqrt{2}}(2\cdot(\ket{000} - \ket{001} - \ket{010} + \ket{011} - \ket{100} + \ket{101} - \ket{110} + \ket{111})) \\ \\
&t_4:& \ket{\psi_4} &= R_{12}\ket{\psi_3} \\
&&     \ket{\psi_4} &= \frac{2}{4\sqrt{2}}\left(
    \begin{aligned}
       &R\ket{00}\ket{0} - R\ket{00}\ket{1} + \\
      -&R\ket{01}\ket{0} + R\ket{01}\ket{1} + \\
      -&R\ket{10}\ket{0} + R\ket{10}\ket{1} + \\
      -&R\ket{11}\ket{0} + R\ket{11}\ket{1} \\
    \end{aligned}
\right) \\
\end{align*}

\begin{align*}
&&     \ket{\psi_4} &= \frac{1}{2\sqrt{2}}(\ket{000} - \ket{001} + \ket{010} - \ket{011} + \ket{100} - \ket{101} + \ket{110} - \ket{111}) \\ 
&&     \ket{\psi_4} &= \frac{1}{2\sqrt{2}}(\ket{0}(\ket{00} - \ket{01} + \ket{10} - \ket{11}) + \ket{1}(\ket{00} - \ket{01} + \ket{10} - \ket{11})) \\
&&     \ket{\psi_4} &= \frac{1}{2\sqrt{2}}(\ket{0} + \ket{1})(\ket{0}(\ket{0} - \ket{1}) + \ket{1}(\ket{0} - \ket{1})) \\
&&     \ket{\psi_4} &= \frac{1}{2\sqrt{2}}(\ket{0} + \ket{1})(\ket{0} + \ket{1})(\ket{0} - \ket{1}) = \ket{+}\ket{+}\ket{-} \\ \\
&t_5:& \ket{\psi_5} &= H_1H_2\ket{\psi_4} \\
&&     \ket{\psi_5} &= \frac{1}{2}(H\ket{0} + H\ket{1})(H\ket{0} + H\ket{1})\ket{-} \\
&t_5:& \ket{\psi_5} &= \frac{1}{4}(\ket{0} + \ket{1} + \ket{0} - \ket{1})(\ket{0} + \ket{1} + \ket{0} - \ket{1})\ket{-} \\
&&     \ket{\psi_5} &= \frac{1}{4}(2\cdot\ket{0})(2\cdot\ket{0})\ket{-} \\
&&     \ket{\psi_5} &= \frac{4}{4}\ket{0}\ket{0}\ket{-} \\
&&     \ket{\psi_5} &= \ket{00}\ket{-} \\
\end{align*}

\subsection*{(b)}
\begin{flushleft}
    Étant donné $\ket{\psi_5} = \ket{00}\ket{-}$, on peut déduire que : \\
\end{flushleft}
\begin{align*}
    prob(x_0 = 0) = prob(\ket{\psi_5} = \ket{000}) + prob(\ket{001}) = (\frac{1}{\sqrt{2}})^2 + (\frac{1}{\sqrt{2}})^2 = 100\% \\
    prob(x_1 = 0) = prob(\ket{\psi_5} = \ket{000}) + prob(\ket{001}) = (\frac{1}{\sqrt{2}})^2 + (\frac{1}{\sqrt{2}})^2 = 100\%\\
\end{align*}

\begin{flushleft}
    Donc, la probabilité que $x = x_1x_0$ est : \\
\end{flushleft}
    \begin{align*} prob(x_1 = 0)\cdot prob(x_0 = 0) = 1 * 1 = 100\% \end{align*}
\newpage
\section*{Question 3}
\subsection*{(a)}
Voici les états avec les calculs des différent $\ket{\psi_1}$ pour $xz = 11$

\begin{align*}
&t_0:&  \ket{\psi_0} &= \ket{1100} \\ \\ 
&t_1:&  \ket{\psi_1} &= H_3\ket{\psi_0} \\
&&      \ket{\psi_1} &= \ket{11}H\ket{0}\ket{0} \\
&&      \ket{\psi_1} &= \ket{11}\frac{1}{\sqrt{2}}(\ket{0} + \ket{1})\ket{0} \\
&&      \ket{\psi_1} &= \frac{1}{\sqrt{2}}(\ket{1100} + \ket{1110}) \\ \\ 
&t_2:&  \ket{\psi_2} &= CNOT_{3\rightarrow4}\ket{\psi_1} \\
&&      \ket{\psi_2} &= \frac{1}{\sqrt{2}}\ket{11}(CNOT_{3\rightarrow4}\ket{00} + CNOT_{3\rightarrow4}\ket{10}) \\
&&      \ket{\psi_2} &= \frac{1}{\sqrt{2}}\ket{11}(\ket{0 \quad 0 \oplus 0} + \ket{1 \quad 1 \oplus 0}) \\
&&      \ket{\psi_2} &= \frac{1}{\sqrt{2}}\ket{11}(\ket{00} + \ket{11}) \\
&&      \ket{\psi_2} &= \frac{1}{\sqrt{2}}(\ket{1100} + \ket{1111}) \\ \\
&t_3:&  \ket{\psi_3} &= X_{2\rightarrow3}\ket{\psi_2} \\
&&      \ket{\psi_3} &= \frac{1}{\sqrt{2}}((X_{2\rightarrow3}\ket{1100}) + (X_{2\rightarrow3}\ket{1111})) \\
&&      \ket{\psi_3} &= \frac{1}{\sqrt{2}}(\ket{1110} + \ket{1101}) \\ \\
&t_4:&  \ket{\psi_4} &= Z_{1\rightarrow3}\ket{\psi_3} \\
&&      \ket{\psi_4} &= \frac{1}{\sqrt{2}}(Z_{1\rightarrow3}\ket{1110} + Z_{1\rightarrow3}\ket{1101}) \\
&&      \ket{\psi_4} &= \frac{1}{\sqrt{2}}(-\ket{1110} + \ket{1101}) = \frac{1}{\sqrt{2}}(\ket{1101} - \ket{1110}) \\ \\
\end{align*}

\begin{align*}
&t_5:&  \ket{\psi_5} &= CNOT_{3\rightarrow4}\ket{\psi_4} \\
&&      \ket{\psi_5} &= \frac{1}{\sqrt{2}}(\ket{11}CNOT_{3\rightarrow4}\ket{01} - \ket{11}CNOT_{3\rightarrow4}\ket{10}) \\
&&      \ket{\psi_5} &= \frac{1}{\sqrt{2}}(\ket{11}\ket{0 \quad 0 \oplus 1} - \ket{11}\ket{1 \quad 1 \oplus 0}) \\
&&      \ket{\psi_5} &= \frac{1}{\sqrt{2}}(\ket{1101} - \ket{1111}) \\ \\ 
&t_6:&  \ket{\psi_6} &= H_3\ket{\psi_5} \\
&&      \ket{\psi_6} &= \frac{1}{\sqrt{2}}((\ket{11}H\ket{0}\ket{1}) - (\ket{11}H\ket{1}\ket{1})) \\
&&      \ket{\psi_6} &= \frac{1}{\sqrt{2}}((\ket{11}\frac{1}{\sqrt{2}}(\ket{0} + \ket{1})\ket{1}) - (\ket{11}\frac{1}{\sqrt{2}}(\ket{0} - \ket{1})\ket{1})) \\
&&      \ket{\psi_6} &= \frac{1}{2}(\ket{1101} + \ket{1111} - \ket{1101} + \ket{1111}) \\
&&      \ket{\psi_6} &= \frac{1}{2}(2\cdot\ket{1111}) = \ket{1111}
\end{align*}
\hspace*{\fill}
\subsection*{(b)}
\begin{flushleft}
    Étant donné $\ket{\psi_6} = \ket{1111}$ et $zx = 11$, on peut déduire que : \\
\end{flushleft}
\begin{align*}
    prob(\hat{z} = 1) = prob{\ket{1111}} = 1^2 = 100\% \\
    prob(\hat{x} = 1) = prob{\ket{1111}} = 1^2 = 100\% \\
\end{align*}
\begin{flushleft}
    Donc, la probabilité que $\hat{z}\hat{x} = zx$ est : \\
\end{flushleft}
    \begin{align*} prob(x_1)\cdot prob(x_0) = 1 \cdot 1 = 100\% \end{align*}
\hspace*{\fill}
\subsection*{(c)}
\begin{align*}
&t_0:&  \ket{\psi_0} &= \ket{zx00} \\ \\
&t_1:&  \ket{\psi_1} &= H_3\ket{zx00} \\
&&      \ket{\psi_1} &= \frac{1}{\sqrt{2}}(\ket{zx00} + \ket{zx10}) \\ \\
&t_2:&  \ket{\psi_2} &= CNOT_{3\rightarrow4}\ket{\psi_1} \\
&&      \ket{\psi_2} &= \frac{1}{\sqrt{2}}(\ket{zx00} + \ket{zx11}) \\ \\
&t_3:&  \ket{\psi_3} &= X_{2\rightarrow3}\ket{\psi_2} \\
&&      \ket{\psi_3} &= \frac{1}{\sqrt{2}}(\ket{zx \quad x \oplus 0 \quad 0} + \ket{zx \quad x \oplus 1 \quad 1}) \\
&&      \ket{\psi_3} &= \frac{1}{\sqrt{2}}(\ket{zxx0} + \ket{zx\overline{x}1}) \\ \\
&t_4:&  \ket{\psi_4} &= Z_{1\rightarrow3}\ket{\psi_3} \\
&&      \ket{\psi_4} &= \frac{1}{\sqrt{2}}((-1)^{z \wedge x}\ket{zxx0} + (-1)^{z \wedge \overline{x}}\ket{zx\overline{x}1}) \\ \\
&t_5:&  \ket{\psi_5} &= CNOT_{3 \rightarrow 4}\ket{\psi_4} \\
&&      \ket{\psi_5} &= \frac{1}{\sqrt{2}}((-1)^{z \wedge x}\ket{zxx \quad x \oplus 0} + (-1)^{z \wedge \overline{x}}\ket{zx\overline{x} \quad \overline{x} \oplus 1}) \\
&&      \ket{\psi_5} &= \frac{1}{\sqrt{2}}((-1)^{z \wedge x}\ket{zxxx} + (-1)^{z \wedge \overline{x}}\ket{zx\overline{x}x}) \\ \\
&t_6:&  \ket{\psi_6} &= H_3\ket{\psi_5} \\
&&      \ket{\psi_6} &= \frac{1}{2}\ket{zx}((-1)^{z \wedge x}(\ket{0x} + (-1)^x \ket{1x}) + (-1)^{z \wedge \overline{x}}(\ket{0x} + (-1)^{\overline{x}} \ket{1x}))
\end{align*}
\newpage
\begin{flushleft}
    Évaluons la fonction pour toutes les valeurs de $z$ et $x$ :
\end{flushleft}
\begin{align*}
    &\ket{zx} = \ket{00} : \\ \\
    \ket{\psi_6} &= \frac{1}{2}\ket{00}((-1)^{0 \wedge 0}(\ket{00} + (-1)^0 \ket{10}) + (-1)^{0 \wedge 1}(\ket{00} + (-1)^1(\ket{10}))) \\
    \ket{\psi_6} &= \frac{1}{2}\ket{00}(\ket{00} + \ket{10} + \ket{00} - \ket{10}) = \ket{0000} \\ \\
    &\ket{zx} = \ket{01} : \\ \\
    \ket{\psi_6} &= \frac{1}{2}\ket{01}((-1)^{0 \wedge 1}(\ket{01} + (-1)^1 \ket{11}) + (-1)^{0 \wedge 0}(\ket{01} + (-1)^0(\ket{11}))) \\
    \ket{\psi_6} &= \frac{1}{2}\ket{01}(\ket{01} - \ket{11} + \ket{01} + \ket{11}) = \ket{0101} \\ \\
    &\ket{zx} = \ket{10} : \\ \\
    \ket{\psi_6} &= \frac{1}{2}\ket{10}((-1)^{1 \wedge 0}(\ket{00} + (-1)^0 \ket{10}) + (-1)^{1 \wedge 1}(\ket{00} + (-1)^1(\ket{10}))) \\
    \ket{\psi_6} &= \frac{1}{2}\ket{10}(\ket{00} + \ket{10} - \ket{00} + \ket{10}) = \ket{1010} \\ \\
    &\ket{zx} = \ket{11} : \\ \\
    \ket{\psi_6} &= \frac{1}{2}\ket{11}((-1)^{1 \wedge 1}(\ket{01} + (-1)^1 \ket{11}) + (-1)^{1 \wedge 0}(\ket{01} + (-1)^0(\ket{11}))) \\
    \ket{\psi_6} &= \frac{1}{2}\ket{11}(-\ket{01} + \ket{11} + \ket{01} + \ket{11}) = \ket{1111} \\ \\
\end{align*}
\begin{flushleft}
    On peut donc voir que le circuit copie dans le $3^e$ et $4^e$ qubit les valeurs de z et x, peut importe leur valeur.
\end{flushleft}
\newpage
\subsection*{(d)}
\begin{flushleft}
    Comme vu à l'énoncé précédent, le circuit "copie" les valeur de $\ket{zx}$ en entrée sur le $3^e$ et $4^e$ qubit.
    La première étape est de partager entre Alice et Bob un état de Bell. La moitié de Alice est celle où l'on effectue
    l'opération XZ, c'est-à-dire sur le $3^e$ qubit sur le schéma. L'information de $\ket{x}$ est directement sur la valeur du $3^e$
    qubit (que l'on copie ensuite sur le $4^e$ qubit) et l'information de $\ket{z}$ est présent dans la phase. Effectuer les portes
    CNOT et H nous permet donc de "retrouver" ces informations à partir du qubit envoyé par Alice.
    sur $\hat{z}\hat{x}$
\end{flushleft}

\section*{Question 4:}
\subsection*{(a)}
La fonction $f_x$ que l'on pourrait donner en entrée au problème \textbf{PF} pour que ça sortie $r$
soit la période $r$ de la suite \{$x^0 = 1\; mod\; N, x^1, x^2, ...\;,$ \\ $x^r = 1\; mod\; N, x^{r+1} = x^1\; mod\; N, ...$\} est
la suivante: \\
\begin{align*} f_x(a) = x^a\; mod\; N \end{align*} \\
Il s'agit de cette fonction, parce qu'il est possible de construire la suite de la manière suivante:
\begin{align*} \{
    f_x(0), f_x(1), f_x(2), ...
    \}
\end{align*} \\
Ce qui veut dire que la période $r$ que \textbf{PF} trouverait pour la fonction $f_x$ serait la même que
la période $r$ que la suite, puisque $f_x$ génère la suite et que \textbf{PF} sert à trouver la période d'une
suite générée par une fonction. On peut appliquer la fonction $f_x$ à \textbf{PF}, parce que celle-ci est périodique
et injective. Celle-ci est injective, parce que l'opération de puissance et celle de modulo sont des opérations de ce style :
\begin{align*} \mathbb{N} \mapsto \mathbb{N} \end{align*} \\
ce qui veut dire que $f_x:\mathbb{N}\mapsto \mathbb{N}$ \\ \\
Aussi, il est possible de dire que la fonction est périodique, puisque, étant donné que notre fonction donne:
\begin{align*} f_x(r) : x^r = 1\; mod\; N \end{align*}
\newpage
\begin{flushleft}
Il est donc possible d'exprimer le résultat pour un nombre $n = c\cdot r + k,\; o\grave{u}\; x >= 1, k \in \{1, ..., N- 1\}$
\end{flushleft}
\begin{align*}
    f_x(n) &= x^n mod N \\
    f_x(n) &= x^{c \cdot r + k} \\
    f_x(n) &= (x^r)^k\cdot x^k \\
    f_x(n) &= 1^k \cdot x^k \\
    f_x(n) &= x^k \\
\end{align*}
ce qui veut dire que: \\ \\
$\{f_x(0),...,f_x(r)\} = \{f_x(r+1), ..., f_x(2\cdot r)\} = \{f_x(c\cdot r + 1), ..., f_x(c\cdot r)\}$ \\ \\
Est donc, la fonction se est donc périodique pour tout $c \in \mathbb{N}$
\end{document}